\section{The adjacent possible and its carriers}
\label{sec:ap}

Of the native formalisms, the adjacent possible is the closest to ICP and the
one in which ICP has the clearest opportunity to contribute. This section sets
out what the formal theory already establishes, identifies the step it leaves
unspecified, and states a prediction that separates a carrier-based account
from the existing models.

\subsection{What is formalized}

Kauffman defined the chemical adjacent possible as the molecular species absent
from the actual but synthesizable from it in one reaction step
\citep[p.~142]{Kauffman2000Investigations}, and later argued in this journal
that enabling constraints open unprestatable possibilities in biological
evolution \citep{Kauffman2014Prolegomenon,PatraKauffman2025NoLaw}. Tria,
Loreto, Servedio, and Strogatz gave the idea its first minimal mathematical
form \citep{TriaEtAl2014Novelties}. In their urn model with triggering, an
element drawn from an urn is returned together with $\rho$ copies
(reinforcement); if and only if the element is new, $\nu+1$ brand-new elements
are also added, so that $\nu+1$ is the size of the adjacent possible opened by
that novelty. The balance of the two processes fixes the statistics: the number
of distinct elements grows as $D(N)\sim N^{\nu/\rho}$ for $\nu<\rho$, as $N\log
N$ for $\nu=\rho$, and linearly for $\nu>\rho$, with Zipf's law $f(R)\sim
R^{-\rho/\nu}$. A semantic-triggering variant reproduces the clustering of
related novelties observed in Wikipedia edits, music listening, annotation tags,
and texts \citep{TriaEtAl2014Novelties}.

The framework has since been extended in several directions. Explorations of
the adjacent possible in social contact space reproduce the rate at which
individuals acquire new acquaintances and the degree, clustering, and community
structure of real networks \citep{UbaldiEtAl2021SocialNetworks}. Waves of
novelties in which newcomers are not suppressed by established elements have
been quantified across several systems \citep{MonechiEtAl2017Waves}. A
time-dependent urn with triggering, with reinforcement and triggering
parameters that vary in time, reproduces the logarithmic novelty growth
observed in gene mutations and patent codes and has a critical region in which
Heaps' and Zipf's laws coexist \citep{BellinaEtAl2025TimeDependentUrn}.
Edge-reinforced walks place discovery on an explicit concept network
\citep{IacopiniEtAl2018Innovation}, and higher-order novelty counts first
occurrences of combinations, distinguishing sequences with similar first-order
discovery rates \citep{DiBonaEtAl2025HigherOrder}. Combinatorial TAP models
generate long plateaus followed by explosive growth
\citep{SteelEtAl2020Combinatorial}, while long-run product data show
resource-constrained rather than unrestricted combinatorial discovery
\citep{Taalbi2026AdjacentPossible}. Fari\~na and Vallejo give a geometric
formulation, a biological cosmological constant proportional to the squared
fractional expansion rate of Fisher-information volume
\citep{FarinaVallejo2026LambdaB}.

ICP therefore cannot claim history-dependent accessibility, triggered support
growth, correlated novelty, a plateau-to-explosion transition, the Heaps and
Zipf regularities, or a geometric measure of expansion. All of these are prior
art.

\subsection{What is left unspecified}

In each of these models the opening step is postulated. The triggering
parameter $\nu$, or its time-dependent generalization, is fitted to data, and
the elements it adds are featureless new colours or nodes; what in the system
produces them lies outside the model. Two further features follow from this
construction. First, the basic urn never removes elements, so possibilities
once opened stay open indefinitely: in the ledger of
Eq.~\eqref{eq:possibility-flux-ledger} only the enabled set $E$ is populated,
and the lost set $L$ is empty by construction. Second, the opened
possibilities depend on the history only through the count and order of past
novelties, not through any structure that remains physically present.

Kauffman's enabling constraints name the missing element without making it
operational. This is the step ICP's carrier criterion is designed for. In its
terms, an account of how the adjacent possible expands must identify an
embodied carrier $m$, formed by the novelty event, whose presence makes the
frontier $\mathcal J^{\pi,\ell}_{t,\tau}$ of Eq.~\eqref{eq:adjacent-frontier}
larger than it is under a matched reference $m_0$. The triggering size then
becomes an intervention--response quantity,
\begin{equation}
 \nu_{\mathrm{C}}(m)
 :=\bigl|\mathcal J^{\pi,\ell}_{t,\tau}(m)\bigr|
  -\bigl|\mathcal J^{\pi,\ell}_{t,\tau}(m_0)\bigr|,
 \label{eq:carrier-triggering}
\end{equation}
measured by ablating, restoring, or replacing the carrier rather than fitted to
a sequence, and reported together with the identities of the enabled and lost
types. ICP does not supply the mechanism that makes $\nu_{\mathrm{C}}$ positive
in a given system; that is chemistry, biophysics, or sociology. It supplies the
test that would identify the mechanism, and the accounting under which the
identification can fail.

\subsection{A discriminating prediction}

The two accounts differ on one point that can be tested. In an urn model,
possibilities opened by a novelty persist whatever happens afterwards. On the
carrier account, an opened possibility persists only while its carrier, or a
role-equivalent successor, remains realized:
\begin{quote}
\textbf{Carrier dependence of the adjacent possible.} If the carrier formed by
a novelty event is removed after the event, without a successor, the
possibilities it opened leave the frontier, and they return if the carrier is
restored.
\end{quote}
A pure triggering model predicts no change in the frontier under this
intervention, because the history of novelties is unchanged. A complete
mechanistic model of a particular system may of course predict the same as the
carrier account; the prediction distinguishes ICP from the phenomenological
models of the adjacent possible, not from complete-state dynamics.

Molecular evolution already provides a case in which the prediction holds.
During the evolution of cortisol specificity in the ancestral glucocorticoid
receptor, permissive substitutions were functionally neutral in isolation but
had to be present for the large-effect function-switching mutations to be
tolerated \citep{HarmsThornton2014Contingency}. In the terms above, the
permissive substitutions are the carrier, the set of tolerable
function-switching mutations is the enabled set $E$, and the dependence of
tolerance on their presence is the carrier dependence of that frontier. The
same study found that permissive genotypes were extremely rare: a library of
thousands of variants of the ancestral protein recovered the historical
permissive substitutions but no alternatives, because a permissive mutation
must stabilize specific local structure, preserve the energetic balance between
conformations, and be compatible with both ancestral and derived structures
\citep{HarmsThornton2014Contingency}. Contingency in which few paths are open
also appears in the exhaustive reconstruction of TEM $\beta$-lactamase
trajectories \citep{WeinreichEtAl2006MutationalPaths}. These are cases in which
the opener of an adjacent possible has been identified physically and in which
its ablation is experimentally meaningful.

The pattern also connects to the crossover of
Eq.~\eqref{eq:constraint-enablement-crossover}. Permissive substitutions
restrict nothing measurable at the level of immediate receptor function, yet
they enlarge the set of later functional changes; a carrier that prunes
microstates while enlarging the role frontier is the general form of the same
observation.

\subsection{From statistics to carriers}

The urn models explain the statistics of discovery from a small number of
parameters. The carrier account asks what sets those parameters. If it is
correct, the triggering and reinforcement rates of a system should be
predictable from measured properties of its carriers: how often novelty events
form a carrier, how large a frontier each carrier opens, and how long carriers
survive or are succeeded. A time-dependent triggering rate
\citep{BellinaEtAl2025TimeDependentUrn} would then be read as a changing
population of realized carriers. Making that link quantitative in one system,
by predicting $\nu$ and its time dependence from carrier measurements before
fitting the sequence, is the most direct way in which ICP could become an
explanation of the adjacent possible rather than a redescription of it.
